\chapter{Infinite divisibility, necklaces and solved models}\label{chap:rzid}\chaptermark{Infinite divisibility and necklaces}

Section~\ref{ch:rz3} proves with no hypothesis the positivity of the canonical function $w$ of Theorem~\ref{rz:thm:w}, which the proof of Theorem~\ref{rz:thm:floor} used, so that $\varphi$ itself is infinitely divisible, and places $\varphi$ in the Bondesson hierarchy: infinitely divisible, but not self decomposable, not in the Bondesson class and not a generalised gamma convolution. It then identifies the Vere-Jones weights as necklace polynomials, extends the L\'evy representation across the critical strip, and compares the solved models with the theory of P\'olya and Cardon. The truncations of this lattice are studied in Volume II.

The coordinates of Chapter~\ref{chap:rzeta} (Section~\ref{rz:sec:intro}) hold here too: in the variable $s$ of the clock a zero $\frac12+i\gamma$ sits at $s=-\gamma^2$ on the cut.
